\documentclass[10pt,a4paper]{amsart}
\usepackage[margin=19mm]{geometry}
\usepackage{amsmath,amssymb,mathtools}
\usepackage[scr=rsfs,cal=euler]{mathalfa}
\usepackage{xfrac}
\usepackage[unicode,hidelinks,bookmarks=false]{hyperref}
\newcommand{\ie}{i.e.\ }
\newcommand{\cf}{cf.\ }
\newcommand{\resp}{respectively\ }
\newcommand{\Subsection}{\subsection}
\providecommand{\nobreakdash}{\nobreak\hbox{-}}
\setlength{\emergencystretch}{2em}
\multlinegap0pt
\usepackage{amssymb}
\usepackage{enumerate}
\usepackage{float}
\usepackage{xcolor}
\newtheorem{theorem}{Theorem}
\newtheorem{lemma}{Lemma}
\title{$k$-Lucas Annulus for Polynomial Zeros}
\author{Herbert Batte}
\date{Source: arXiv:2608.26452v1 (CC BY 4.0)}
\begin{document}
\maketitle

\noindent\emph{Source and licence.} $k$-Lucas Annulus for Polynomial Zeros. Version arXiv:2608.26452v1 (CC BY 4.0), distributed by arXiv under the Creative Commons Attribution 4.0 International licence, http://creativecommons.org/licenses/by/4.0/, as recorded in the arXiv licence metadata for this exact version and shown on the abstract page https://arxiv.org/abs/2608.26452v1. Full licence notice: \texttt{licenses/arxiv-2608.26452-CC-BY-4.0.txt}.\par\smallskip

\noindent\textit{Editorial scope.} Counted unit: the article's theorem thm:main (source0.tex lines 84--95). The article's complete original proof of it is delivered with it (source0.tex lines 209--244, one \texttt{proof} environment of 36 lines, which the article places away from the statement). No mathematics is written, repaired or re-derived: the statement and the proof are the article's own lines, in the article's own notation, and no retained line is substituted or re-flowed.  Retained uncounted as the source-local context the proof needs: lines 53--53; lines 158--166.  Nothing was omitted from any retained span, so the package carries no omission record.  No retained line contains a citation, so the wrapper carries no bibliography and no bibliography entry.  No retained line contains a figure environment, an image-inclusion command or drawing code, so no image is delivered and the package declares no asset dependency.  Editorial formatting only, in addition: an AMS \texttt{amsart} preamble that restates the article's own declarations and macros used by the retained lines, namely the environments \texttt{theorem}, \texttt{lemma} on separate counters, together with the standard packages those lines need.  The retained environments print in this document as Theorem 1, Lemma 1 in order of appearance, whereas the article prints the same items with the numbering of its own full document.  The complete native source of the article as distributed in the arXiv e-print for this exact version is bundled unchanged as \texttt{sources/208633/source0.tex}.
\subsection{Background}\label{sec:background}

\begin{lemma}[Key identity]
	\label{lem:key}
Let $\alpha,\beta$ be as in Section~\ref{sec:background}, and let $m\ge1$ be an integer. Then
	\[
	L_{k,m}-\beta L_{k,m-1}=\alpha^{m-1}(\alpha-\beta),
	\qquad
	L_{k,m}-\alpha L_{k,m-1}=-\beta^{m-1}(\alpha-\beta).
	\]
\end{lemma}

\begin{theorem}[$k$-Lucas -- $k$-Fibonacci binomial identity]
	\label{thm:main}
	Let $k>0$ be a real number, and let $m,n$ be integers with $m\ge 1$ and $n\ge1$. Then
 \begin{align*}
	\sum_{\ell=0}^{n}\binom{n}{\ell}(L_{k,m-1})^{n-\ell}(L_{k,m})^{\ell}L_{k,\ell}
	=
	\begin{cases}
		(k^2+4)^{n/2}\,L_{k,mn}, & \text{if } n \text{ is even},\\[6pt]
		(k^2+4)^{(n+1)/2}\,F_{k,mn}, & \text{if } n \text{ is odd}.
	\end{cases}
 \end{align*}
\end{theorem}

	\begin{proof}[Proof of Theorem \ref{thm:main}]
		Write $S=\sum_{\ell=0}^n\binom{n}{\ell}(L_{k,m-1})^{n-\ell}(L_{k,m})^{\ell}L_{k,\ell}$. Using $L_{k,\ell}=\alpha^{\ell}+\beta^{\ell}$,
		\[
		S=\sum_{\ell=0}^n\binom{n}{\ell}(L_{k,m-1})^{n-\ell}(L_{k,m})^{\ell}\alpha^{\ell}
		+\sum_{\ell=0}^n\binom{n}{\ell}(L_{k,m-1})^{n-\ell}(L_{k,m})^{\ell}\beta^{\ell}
		=:S_1+S_2.
		\]
		By the binomial theorem,
		\[
		S_1=(L_{k,m-1}+\alpha L_{k,m})^n,\qquad S_2=(L_{k,m-1}+\beta L_{k,m})^n.
		\]
		Since $\alpha,\beta$ are the roots of $x^2-kx-1=0$, we have $\alpha\beta=-1$. Hence, by Lemma \ref{lem:key},
		\[
		L_{k,m-1}+\alpha L_{k,m}=\alpha\bigl(L_{k,m}-\beta L_{k,m-1}\bigr)=\alpha\cdot\alpha^{m-1}(\alpha-\beta)=\alpha^{m}(\alpha-\beta),
		\]
		\[
		L_{k,m-1}+\beta L_{k,m}=\beta\bigl(L_{k,m}-\alpha L_{k,m-1}\bigr)=\beta\cdot\bigl(-\beta^{m-1}(\alpha-\beta)\bigr)=-\beta^{m}(\alpha-\beta).
		\]
		Therefore
		\[
		S_1=\alpha^{mn}(\alpha-\beta)^n,\qquad S_2=(-1)^n\beta^{mn}(\alpha-\beta)^n,
		\]
		so
		\[
		S=(\alpha-\beta)^n\bigl(\alpha^{mn}+(-1)^n\beta^{mn}\bigr).
		\]
		Since $\alpha-\beta=\sqrt{k^2+4}$, we have $(\alpha-\beta)^n=(k^2+4)^{n/2}$.
		
		If $n$ is even, then $\alpha^{mn}+(-1)^n\beta^{mn}=\alpha^{mn}+\beta^{mn}=L_{k,mn}$, and $S=(k^2+4)^{n/2}L_{k,mn}$.
		
		If $n$ is odd, then $\alpha^{mn}+(-1)^n\beta^{mn}=\alpha^{mn}-\beta^{mn}=(\alpha-\beta)F_{k,mn}$, using the Binet formula for $F_{k,mn}$, and so
		\[
		S=(\alpha-\beta)^n\cdot(\alpha-\beta)F_{k,mn}=(\alpha-\beta)^{n+1}F_{k,mn}=(k^2+4)^{(n+1)/2}F_{k,mn}.
		\]
		This establishes both cases.
	\end{proof}

\end{document}
